\documentclass{article}
\usepackage[utf8]{inputenc}
\usepackage{graphicx}
\usepackage{geometry}
\usepackage{hyperref}
\hypersetup{
    colorlinks,
    citecolor=black,
    filecolor=black,
    linkcolor=black,
    urlcolor=black
}

 \geometry{
 a4paper,
 left=30mm,
 right=30mm,
 top=30mm,
 }

\graphicspath{ {Images/} }

%----------------------------------------------------------------------------------------
%	TITLE PAGE
%----------------------------------------------------------------------------------------

\newcommand*{\titleGP}{\begingroup
		\begin{figure}[t]
			\centering
			\includegraphics[width=350px]{UP_Logo.PNG}
		\end{figure}
\centering 
\vspace*{\baselineskip}

\rule{\textwidth}{1.6pt}\vspace*{-\baselineskip}\vspace*{2pt}
\rule{\textwidth}{0.4pt}\\[\baselineskip]

{\LARGE NavUP\\ [0.3\baselineskip] Longsword Testing Report } \\ [0.2\baselineskip]
\rule{\textwidth}{0.4pt}\vspace*{-\baselineskip}\vspace{3.2pt}
\rule{\textwidth}{1.6pt}\\[\baselineskip] %

% \scshape %
% A concise specification on the functional requirements  \\
% and use cases of NavUP \\[\baselineskip]

% \vspace*{2\baselineskip}

Compiled By \\[\baselineskip]
{\Large Lucian Sargeant - u15225560 \\ Ritesh Doolabh - u15075754 \\ Peter Boxall -  u14056136 \\ Claude Greeff - u13153740\\ Harris Leshaba - u15312144 \\ Hristian Vitrychenko - u15006442\par}

\bigskip
\bigskip

 	GitHub Repository:  
 	\href{https://github.com/Chris19951225/COS-301-Longsword-Data-Streaming}{COS 301 Team Longsword Data GitHub Repository(Phase 4)}




 

\vfill


{\scshape 2017} \\[0.3\baselineskip]
{\large TEAM LONGSWORD (DATA)}\par

\endgroup}

\begin{document}

\titleGP
\newpage
\tableofcontents

\newpage
\section{Introduction}
\begin{flushleft}
For this phase we will be testing the Data module of the BroadSword Team. We have split the testing phase according to Functional Requirements, Non-Functional Requirements and Use Cases. 
\end{flushleft}

\begin{flushleft}
Their code was primarily coded in Python and used a NSQ message processing system.
We will be testing the various cases and giving a brief description of how we tested followed by an explanation of the mark that was given to them.
\end{flushleft}


\section{Service Contracts}

\subsection{Retrieving and passing device MAC address.}
\begin{flushleft}MARK: 10\end{flushleft}
The access module sends a location request to the data module via the NSQ server,the access module has to subscribes to a topic.The main entry point of the system, query\_resolver.py, sends the mac address to the aruba server in a query to get the location of the device.
In query\_resolver.py, the handler function recieves a request, it then passes the mac address to the searcher function-which will create a query to get the location for the device.
This functional requirement was fullfilled by the Broadsword team.
\begin{figure}[t]
	\centering
	\includegraphics{handler.PNG}
\end{figure}
\subsection{Logging in and maintaining a session with Aruba ALE.}

\subsection{Processing the request and retrieving location.}

\subsection{Returning a location to the source of the request.}


\section{Non-Functional Requirements}

\subsection{Level of concurrency of the task.}

\subsection{Performance of the request processing (time taken to receive a response).}

\subsection{Maintainability and modularity of the code and repository.}

\subsection{Integrability and ease of transfer into a final system.}


\section{Use Cases}

\subsection{Upstream communication.}

\subsection{Downstream communication.}

\end{document}
